% Preface and conventions, book-only front matter. Input from aic.tex
% between the table of contents and the chapters.

\setchapterstyle{plain}
\chapter*{Preface}
\addcontentsline{toc}{chapter}{Preface}

These are the lecture notes for \emph{TFE4188 Advanced Integrated
Circuits} at the Norwegian University of Science and Technology, grown
over years of teaching into something that behaves like a book. The
intended reader is a student who has met circuits before and now wants
to design integrated ones: analog on a modern CMOS process, with the
digital world pressing in from every side.

The book has two kinds of chapters. The \emph{refresher} chapters -
physics, fields, diodes, MOSFETs, circuits, OTAs, passives, noise,
logic, layout - rebuild the foundations the later chapters stand on;
read them in order the first time, and raid them afterwards. The
\emph{technical} chapters - ESD, references, filters, converters,
regulators, clocks, oscillators, radio - each take one job an
integrated circuit must do and follow it from principle to schematic.
Every chapter ends with a one-page summary and a reading list.

Everything here runs on an open source flow: xschem, ngspice, magic,
and the cic tools, on the open SKY130 process. Every schematic figure
can be simulated, every simulation rerun from a shell command, and the
source of the book itself - text, figures, build - lives in a public
git repository. If you find an error, the repository is where to say
so.

Parts of this book were written by Claude, Anthropic's AI, on my
direction and under my review; the title page says so, AI-written
chapters carry a byline, and the repository's commit history records
precisely who wrote what.

\bigskip
\noindent Carsten Wulff, Trondheim, 2026

\setchapterstyle{plain}
\chapter*{Conventions}
\addcontentsline{toc}{chapter}{Conventions}

The figures in this book follow one visual language, so colour and
line style carry meaning everywhere:

\begin{itemize}
\item \textbf{Black} is the circuit: wires, devices, and the house
  symbols. One line width throughout, one arrow tip.
\item \textcolor{red}{\textbf{Red}} marks the subject - the device or
  path the figure is about. In the device cartoons red is also the
  hole, and red regions are p-type.
\item \textcolor{blue}{\textbf{Blue}} carries signals and fields. In
  the device cartoons blue is the electron, and blue regions are
  n-type.
\item \textcolor{armygreen}{\textbf{Green}} marks equivalence - two
  drawings of the same thing - and mechanical annotations.
\item {\color[RGB]{60,140,160}\textbf{Teal}} is the gate oxide,
  wherever it appears.
\item \textbf{Dashed or arrowed strokes} are annotation, never wires.
\item \textbf{Wires that cross without a dot are not connected};
  every junction carries a dot.
\end{itemize}

Device sizes quoted as \texttt{24F4F} read as width 24 and length 4
minimum gate lengths, so a schematic ports between processes by
re-reading F. Some numerical estimates are deliberately rough -
engineering hand-waving is a skill this book practices openly, and
the text says so where it happens.
